\begin{apartats}

\apartat[1,25]{0,75}
Escriu les equacions vectorial, paramètriques i contínua de la recta $r$ que passa pel punt $A(2,-1,3)$
i té vector director $\vec v=(1,-2,2)$.

\begin{solucio}
Vectorial: $(x,y,z)=(2,-1,3)+t(1,-2,2)$.
\[
\text{Paramètriques: }\left\{\begin{aligned}x&=2+t\\y&=-1-2t\\z&=3+2t\end{aligned}\right.
\qquad
\text{Contínua: }\frac{x-2}{1}=\frac{y+1}{-2}=\frac{z-3}{2}.
\]
\end{solucio}

\begin{tria}{recta-punts}
\itemtria{dos-punts}{1}{1,25}
Escriu les equacions contínua i implícita de la recta que passa per $B(1,0,2)$ i $C(3,-1,5)$.

\begin{solucio}
Vector director: $\overrightarrow{BC}=(2,-1,3)$. Contínua: $\dfrac{x-1}{2}=\dfrac{y}{-1}=\dfrac{z-2}{3}$.\\
Igualant dues a dues: $-(x-1)=2y$ i $3(x-1)=2(z-2)$. Implícita:
$\left\{\begin{aligned}x+2y-1&=0\\3x-2z+1&=0\end{aligned}\right.$
\end{solucio}

\itemtria{alineats}{1}{1,25}
Troba el valor de $a$ perquè els punts $P(1,2,-1)$, $Q(3,1,1)$ i $R(7,a,5)$ estiguin alineats.

\begin{solucio}
Estan alineats si $\overrightarrow{PQ}=(2,-1,2)$ i $\overrightarrow{PR}=(6,a-2,6)$ són proporcionals:
$\dfrac62=\dfrac{a-2}{-1}=\dfrac62$. La primera i la tercera components ja donen la raó 3, i cal
$a-2=-3$, és a dir, $a=-1$.
\end{solucio}
\end{tria}

\begin{nomesllarg}
\apartat{0,75}
Escriu dos punts de la recta $r$ del primer apartat, diferents de $A$. Pertany a $r$ el punt $S(5,-7,9)$?

\begin{solucio}
Resposta oberta. Amb $t=1$, $(3,-3,5)$, i amb $t=-1$, $(1,1,1)$.\\
$S$ hi pertany si hi ha un mateix $t$ per a les tres coordenades: $2+t=5$ dona $t=3$, i amb $t=3$,
$y=-1-6=-7$ i $z=3+6=9$. Sí que hi pertany.
\end{solucio}
\end{nomesllarg}

\end{apartats}
